% !TeX root = ../../Book4.tex
% 纲要标题：### 20.4 同调代数的进一步发展
% 纲要位置：计划纲要.md，第 3045 行。

\section{同调代数的进一步发展}
\label{b4:sec:2004}


% 纲要内容（仅作写作计划，尚非教材正文）：
%
% * 增强的用途
% * 稳定高阶范畴的动机
% * 几何表示论与导出几何所需的附加语言
%
% ---
%

\subsection{增强的用途}
\label{b4:subsec:2004:01}

增强的用途可以从一个很小的计算看出。设 \(R\) 为含幺环，\(f:C\to D\)、\(f':C'\to D'\)、\(u:C\to C'\)、\(v:D\to D'\) 为复形映射。若方块严格交换，便有锥之间的对角映射 \((y,x)\mapsto(vy,ux)\)。若方块只在同伦范畴中交换，实际给出的条件则是存在次数 \(-1\) 的 \(h:C\to D'\)，使
\[
 vf-f'u=d_{D'}h+hd_C.
\]
此时锥之间的映射为
\begin{equation}\label{b4:eq:cone-map-with-homotopy}
 \operatorname{Cone}(f)\longrightarrow\operatorname{Cone}(f'),
 \qquad (y,x)\longmapsto(vy+hx,ux).
\end{equation}
把两边微分写成 \(\left(\begin{smallmatrix}d_D&f\\0&-d_C\end{smallmatrix}\right)\)，矩阵相乘后，右上角的交换条件恰好就是所给同伦等式。因此锥映射需要 \(h\) 这项具体资料；只知道 \([vf]=[f'u]\) 并没有指定它。

\ProseParagraphBreak
这种依赖确实可能改变所得映射。在域 \(k\) 上，取 \(C=C'=k[-1]\)、\(D=D'=k\)、\(f=f'=0\)、\(u=v=0\)。次数 \(-1\) 的 \(h:C\to D'\) 可以是任意标量。此时两个锥都是集中在零次的 \(k\oplus k\)，\eqref{b4:eq:cone-map-with-homotopy}为 \(\left(\begin{smallmatrix}0&h\\0&0\end{smallmatrix}\right)\)。不同标量给出不同同伦类，因为这些锥没有非零微分，也没有非零同伦。一个在同伦范畴里固定的零方块，仍然允许不同锥映射。

\ProseParagraphBreak
DG 增强保留整个态射复形，从而保留同伦所在的负一次以及同伦之间更高次的比较。它不会神奇地消除选择，而是把选择及其差别放在可以继续计算的对象中。对已有的三角等价，是否存在相容的 DG 提升是额外问题；本章只构造了具体模型中的增强，并使用特定双模复形给出的等价，没有断言任意三角等价都能如此提升。

\subsection{稳定高阶范畴的动机}
\label{b4:subsec:2004:02}

上一小节中，\(h\) 是两个映射之间的同伦；若两个同伦的差又是 \(\partial s\)，其中 \(s\) 的次数为 \(-2\)，便有同伦之间的比较。继续下去，会出现一串彼此相容的高阶资料。普通范畴只记录对象和态射；三角范畴额外记录移位与好三角，但并不逐层保留这些资料。

\ProseParagraphBreak
\term[wendinggaojiefanchou]{稳定高阶范畴}[stable infinity-category]研究这种现象时，把态射集合换为包含路径及高阶同伦的映射空间。这里仅说明名称所对应的目标，不给出该理论的完整定义与公理。稳定性要求一种带零对象、有限极限和有限余极限的高阶范畴结构，并要求拉回方块与推出方块相同；在这种结构中，纤维与余纤维通过移位相联系。仅有普通范畴中的这些词还不足以构成定义，因为极限、相容图式及唯一性都必须在高阶意义下解释。

\ProseParagraphBreak
对复形而言，映射纤维可写成 \(\operatorname{Cone}(f)[-1]\)，映射余纤维可写成 \(\operatorname{Cone}(f)\)。\chapref{b4:ch:17}证明的三角关系已经显示了它们与移位的联系；高阶理论进一步把构造所需的相容同伦一起保存。因此学习它之前，先熟悉锥、导出函子和本章的 Hom 复形，能够明确看见新增结构要解决什么问题。本卷到此只提供动机，不由三角公理推导高阶范畴的存在或唯一性。

\subsection{几何表示论与导出几何简介}
\label{b4:subsec:2004:03}

几何中的一个基本现象，是两个子空间的交会出现额外的同调。设 \(k\) 为域，\(A=k[x]\)、\(M=A/(x)\)。普通张量积只给出 \(M\otimes_A M=k\)。若先用自由消解
\[
 P=(A\xrightarrow{x}A),\qquad P^{-1}=P^0=A,
\]
再张量 \(M\)，得到
\[
 M\otimes_A^{\mathbf L}M\simeq(k\xrightarrow{0}k),
 \qquad H^{-1}=k,\quad H^0=k.
\]
非零负次同调记录了关系 \(x\) 在再次模去 \(x\) 后不再提供独立约束。若保留\cref{b4:exm:dg-koszul-one-element}的乘法，这个计算还得到 \(k[\varepsilon]/(\varepsilon^2)\)、\(|\varepsilon|=-1\)、\(d=0\) 的 DG 代数。导出几何把类似资料用于研究交、变形与模问题，而不只保留零次商环。

\ProseParagraphBreak
这个算例并不已经定义一般的导出概形。要从仿射计算推广到几何对象，还需要层与下降、局部模型之间的粘合，并说明采用哪一种增强；在正特征中，也不能未经比较就把交换 DG 代数当成所有导出交换环的通用模型。第三卷已经介绍普通层与概形，但一般导出粘合、余切复形及其形变意义不在本章建立。

\ProseParagraphBreak
几何表示论则在空间或概形上研究带群作用的层、它们的上同调以及卷积运算，再从这些结构构造和比较表示。这里不仅需要本卷的导出范畴，还需要具体的层理论、等变性与相应的几何方法。第五卷先承担有限群、箭图和 Lie 代数表示的代数基础；本节只是说明本卷的同调构造为何会在这些方向再次出现，不把这些方向列作本卷已经完成的理论。
